\SourceLeaf{102}
\expandafter\def\csname HistoricalNoteText10\endcsname{E. g. amplitudo oscillationum acus primae in experimento primo fuit initio $0^0\ 37'26''$, in fine $0^0\ 28'34''$; in experimento quinto initio $1^0\ 10'21''$, post $177$ oscillationes $0^0\ 45'35''$, post $677$ oscillationes $0^0\ 6'44''$.}
\begin{center}
\begin{tabular}{c||c|c||c}
&\multicolumn{3}{c}{Oscillationes simultaneae}\\
Exp.&\multicolumn{2}{c||}{acus primae}&acus secundae\\\cline{1-4}
&Oneratio&una oscillatio&una oscillatio\\
I&$r=180^{\mathrm{mm}}$&$24''63956$&$17''32191$\\
II&$r=130$&$20{,}77576$&$17{,}32051$\\
III&$r=80$&$17{,}66798$&$17{,}31653$\\
IV&$r=30$&$15{,}80310$&$17{,}30529$\\
V&sine oneratione&$15{,}22990$&$17{,}31107$
\end{tabular}
\end{center}

Tempora observata sunt ad chronometrum, cuius retardatio intra diem temporis medii erat $14''24$; utrumque pondus $p$ erat $103{,}2572$ grammatum; distantiae $r$ in millimetris microscopica praecisione determinatae; duratio unius oscillationis ad minimum ex $100$ oscillationibus (in experimento quinto adeo ex $677$ pro acu prima) conclusa reductionem ad infinite parvam iam accepit: ceterum hae reductiones propter perparvam oscillationum amplitudinem\HistoricalNote{10}, quam in apparatibus nostris salva summa praecisione adhibere licet, insensibiles sunt. Haec tempora oscillationum reducemus, primo ad valorem medium ipsius $TM$, qui inter experimentum quintum locum habuit, adiumento praeceptorum art. praec II.; dein ad valores, qui absque torsione proventuri fuissent, multiplicatione per $\BareRoot{\frac{n+1}{n}}$, ubi $n$ in quatuor primis experimentis $=424{,}8$, in quinto $=597{,}4$ (conf. art. $9$); denique ad tempus solare medium multiplicatione per $\frac{86400}{86385{,}76}$: hoc pacto nanciscimur
\begin{center}
\begin{tabular}{crl}
I.&$24''65717=t'$&pro $r'=180^{\mathrm{mm}}$\\
II.&$20{,}79228=t''$&pro $r''=130^{\mathrm{mm}}$\\
III.&$17{,}68610=t'''$&pro $r'''=80^{\mathrm{mm}}$\\
IV.&$15{,}82958=t''''$&pro $r''''=30^{\mathrm{mm}}$\\
V.&$15.24515=t$&pro acu non onerata.
\end{tabular}
\end{center}

Accipiendo pro unitatibus temporis, distantiae et massae minutum secundum, millimetrum et milligramma, ut sit $p=103257.2$, e combinatione experimenti primi cum quarto deducimus:
\SourceLeaf{103}
\[
TM=179641070,\qquad K+C=4374976000
\]
ac dein ex experimento quinto
\[
K=4230282000,\quad \text{nec non}\quad C=144694000
\]
Si vero cuncta experimenta ad calculum revocare placet, methodus quadratorum minimorum commodissime sequenti modo applicatur. Proficiscimur a valoribus approximatis incognitarum $x,y$ e combinatione experimenti primi et quarti prodeuntibus, denotatisque correctionibus adhuc adiiciendis per $\xi,\eta$, statuimus
\[
\begin{aligned}
x&=88{,}13646+\xi\\
y&=21184{,}85+\eta
\end{aligned}
\]
Hoc pacto valores calculati temporum $t'.\ t'',t''',t''''$ prodeunt per methodos obvias
\[
\begin{aligned}
t'&=24{,}65717-0{,}13988\xi+0{,}00023008\eta\\
t''&=20{,}78731-0{,}11793\xi+0{,}00027291\eta\\
t'''&=17{,}69121-0{,}10036\xi+0{,}00032067\eta\\
t''''&=15{,}82958-0{,}08980\xi+0{,}00035838\eta
\end{aligned}
\]
quorum comparatio cum valoribus observatis secundum methodum quadratorum minimorum tractata suppeditat
\[
\begin{aligned}
\xi&=-0{,}03230,&\quad\eta&=-12{,}38\\
x&=88{,}10416,&y&=21172{,}47
\end{aligned}
\]

Hinc denique prodit
\[
TM=179575250,\qquad K+C=4372419000
\]
ac dein per experimentum primum
\[
K=4228732400,\qquad C=143686600
\]

Ecce comparationem temporum e valoribus correctis quantitatum $x,y$ calculatorum cum observatis:
\begin{center}
\begin{tabular}{c|r|r|r}
Exper.&Tempus calculatum&Tempus observatum&Differentia\\\hline
I&$24''65884$&$24''65717$&$+0''00167$\\
II&$20{,}78774$&$20{,}79228$&$-0{,}00454$\\
III&$17{,}69046$&$17{,}68610$&$+0{,}00436$\\
IV&$15{,}82805$&$15{,}82958$&$-0{,}00153$
\end{tabular}
\end{center}

\SourceLeaf{104}
Longitudinem penduli simplicis Gottingae statuimus $=994^{\mathrm{mm}}126$, unde fit gravitas, per eam unitatem virium acceleratricium, quae calculis praecedentibus subest, mensurata, $=9811{,}63$: quodsi itaque gravitatem ipsam pro unitate accipere malumus, fit $TM=18302{,}29$: hic numerus exprimit multitudinem milligrammatum, quorum pressio, sub actione gravitatis, in vectem, cuius longitudo est millimetrum, aequivalet vi, qua magnetismus terrestris acum illam circa axem verticalem vertere nititur.

\Article{12}
Postquam determinationem producti vis magneticae terrestris horizontalis $T$ in momentum magnetismi acus datae $M$ absolvimus, iam ad alteram disquisitionis partem progredimur, puta ad determinationem quotientis $\frac MT$. Quam assequemur per comparationem actionis istius acus in aliam acum cum actione magnetismi terrestris in eandem, et quidem, uti iam in introductione expositum est, haec vel in statu motus vel in statu aequilibrii observari poterit: utramque methodum frequenter experti sumus; sed quum posterior pluribus rationibus priori longe praeferenda sit, hocce quidem loco disquisitionem ad illam restringemus. praesertim quum prior prorsus simili modo absque difficultate tractari possit.

\Article{13}
Conditiones aequilibrii corporis mobilis, in quod vires quaecunque agunt, per principium motuum virtualium perfacile in formulam unicam contrahuntur, scilicet aggregatum productorum singularum virium per motum infinite parvum puncti, in quod quaelibet agit, in huius directionem proiectum, esse debet tale, ut pro nullo motu virtuali, i. e. cum conditionibus generalibus, quibus motus corporis subiectus est, conciliabili, valorem positivum obtinere possit, adeoque, quatenus motus virtuales in partes oppositas ubique possibiles sunt, ut illud aggregatum, quod per $\mathrm{d}\Omega$ denotabimus, fiat $=0$ pro quolibet motu virtuali.

Corpus mobile, quod hic consideramus, est acus magnetica, cuius punctum $G$ filo torsili superne fixo annexum est. Hoc filum tantummodo impedit, quominus distantia puncti $G$ a termino fili fixo fieri possit maior longitudine fili, ita ut hic quoque, ut in casu corporis perfecte liberi, positio corporis in spatio a sex variabilibus, adeoque eius aequilibrium a sex conditionibus pendeat: sed quum hoc loco problematis solutio tantummodo determinationi quotientis $\frac MT$ inservire \SourceLeaf{105}debeat, sufficit consideratio motus virtualis eius, qui in rotatione circa axem verticalem per $G$ transeuntem consistit, manifestoque talem axem tamquam fixum et solum angulum inter planum verticale, in quo est acus axis magneticus, atque planum meridianum magneticum tamquam variabilem considerare licebit. Hunc angulum a parte meridiani boreali versus orientem numerabimus et per $u$ denotabimus.

\Article{14}
Concipiamus volumen acus mobilis in elementa infinite parva divisum, sintque $x,y,z$ coordinatae elementi indefiniti, atque $e$ elementum magnetismi liberi in ipso contentum. Initium coordinatarum collocamus in rectae verticalis per $G$ transeuntis puncto arbitrario $h$ intra acum; axes coordinatarum $x,y$ sunto horizontales, ille in meridiano magnetico boream versus, hic versus orientem; coordinatam $z$ sursum numeramus. Ita actio magnetismi terrestris in elementum $e$ producit partem ipsius $\mathrm{d}\Omega$ hancce $Te\mathrm{d}x$.

Simili modo dividatur volumen acus secundae fixae in elementa infinite parva, respondeantque elemento indefinito coordinatae $X,Y,Z$, atque quantitas magnetismi liberi $E$; denique sit $r=\BareRoot{((X-x)^2+(Y-y)^2+(Z-z)^2)}$. Hoc pacto actio elementi $E$ in elementum $e$ sistit partem aggregati $\mathrm{d}\Omega$ hanc $\frac{eE\mathrm{d}r}{r^n}$, si potestati $r^n$ distantiae $r$ reciproce proportionalis supponitur.

Denotando per $N$ eum valorem ipsius $u$, qui detorsioni fili respondet, momentum vis torsionis fili per $\theta(N-u)$ exprimi poterit: haec vis ita concipi potest, ac si in diametri horizontalis fili ad punctum $G$ terminum utrumque ageret vis tangentialis $=\frac{\theta(N-u)}{D}$, denotante $D$ hunc diametrum, unde facile perspicitur, hinc prodire partem aggregati $\mathrm{d}\Omega$ hanc $\theta(N-u)\mathrm{d}u$.

Gravitas particularum acus manifesto nihil confert ad aggregatum $\mathrm{d}\Omega$, quum $u$ sit unica variabilis, quapropter habemus
\[
\mathrm{d}\Omega=\Sigma Te\mathrm{d}x+\Sigma\frac{eE\mathrm{d}r}{r^n}+\theta(N-u)\mathrm{d}u
\]
ubi summatio in termino primo refertur ad cuncta elementa $e$, in secundo ad cunctas combinationes singulorum $e$ cum singulis $E$. Patet itaque, conditionem aequilibrii stabilis consistere in eo, ut
\[
\Omega=\Sigma Tex-\Sigma\frac{eE}{(n-1)r^{n-1}}-\tfrac12\theta(N-u)^2
\]
fiat maximum.

\SourceLeaf{106}
\Article{15}
Ad propositum nostrum convenit, experimenta ita semper adornare, ut axis magneticus utriusque acus sit horizontalis, atque utraque acus in eadem fere altitudine: his itaque suppositionibus calculos ulteriores adstringemus.

Referamus coordinatas punctorum primae acus ad axes in hac fixos in puncto $h$ etiamnum se secantes, et quidem sit axis primus in directione axis magnetici, secundus horizontalis primoque ad dextram, tertius verticalis sursum directus: coordinatae elementi $e$ respectu horum axium sint $a,b,c$. Perinde sint $A,B,C$ coordinatae elementi $E$ respectu similium axium in acu secunda fixorum et in puncto $H$ huius acus se secantium: hoc punctum prope medium acus atque in eadem altitudine cum puncto $h$ electum supponimus.

Situs puncti $H$ commodissime quidem per distantiam a puncto $h$ atque directionem rectae iungentis determinaretur, si de uno tantum experimento ageretur: sed quum ad institutum nostrum semper plura experimenta requirantur ad diversas puncti $H$ positiones spectantia, quae quidem omnes sunt in eadem recta, attamen haud necessario in recta per punctum $h$ exacte transeunte, praestat, signa statim ab initio ita adornare, ut systema talium experimentorum ab unica variabili pendeat. Referemus itaque punctum $H$ ad punctum arbitrarium $h'$ in eodem plano horizontali ipsi $h$ propinquum, cuius coordinatae sint $\alpha,\beta,0$, statuemusque distantiam $h'H=R$. angulumque rectae $h'H$ cum meridiano magnetico $=\psi$. Quodsi iam angulum axis magnetici secundae acus cum meridiano magnetico per $U$ denotamus, habebimus
\[
\begin{aligned}
x&=a\cos u-b\sin u\\
y&=a\sin u+b\cos u\\
z&=c\\
X&=\alpha+R\cos\psi+A\cos U-B\sin U\\
Y&=\beta+R\sin\psi+A\sin U+B\cos U\\
Z&=C
\end{aligned}
\]

Ita omnia ad evolutionem aggregati $\Omega$, atque quotientis $\frac{\mathrm{d}\Omega}{\mathrm{d}u}$, qui pro statu aequilibrii evanescere debet, praeparata sunt.

\Article{16}
Primo fit $\Sigma Tex=T\cos u\, .\, \Sigma ae-T\sin u\, .\, \Sigma be=mT\cos u$, si momen%
